\documentclass[11pt]{article}
\usepackage[margin=1in]{geometry}
\usepackage{amsmath,amssymb,amsfonts}
\usepackage{bm}
\usepackage{hyperref}
\usepackage{physics}
\usepackage{enumitem}
\usepackage{graphicx}
\usepackage{xcolor}

\hypersetup{
  colorlinks=true,
  linkcolor=blue!60!black,
  citecolor=blue!60!black,
  urlcolor=blue!60!black
}

\title{Deterministic Path Ontologies in Interferometry:\\
Hidden Determinism (PCC+QHP) vs.\ de Broglie--Bohm Trajectories\\[4pt]}
\author{Citizen Gardens \\ The Foundation of Multiplicity}
\date{May 4, 2026}

\begin{document}
\maketitle

\begin{abstract}
This report formalizes and unifies a set of developments concerning deterministic trajectory-based interpretations of quantum interferometry, with emphasis on the Hidden Deterministic Interpretation (HDI) of Gurappa and its Phase Consistency Criterion (PCC) and Quantum Hamilton's Principle (QHP) on one side, and de Broglie--Bohm (dBB) theory on the other.\,[file:1][web:52]
We provide (i) an executive-level summary, (ii) a comprehensive mathematical description of PCC+QHP in a path-integral language, (iii) an explicit identification of the core incompatibility between PCC+QHP and Bohmian trajectories in double-slit and Mach--Zehnder interferometer (MZI) setups, (iv) a unifying numerical architecture based on a 2D time-dependent Schr\"odinger solver with split-step Fourier propagation, and (v) concrete code templates and design patterns for interactive simulators that visually expose these ontological differences while keeping the underlying wave evolution fixed.\,[file:1][web:67][web:79][web:96]
The document is written as a defensive publication: all constructions, simulation patterns, and specific algorithmic combinations are disclosed as prior art.
\end{abstract}
\newpage
\tableofcontents

\section{Executive Summary}

\subsection{Conceptual core}

The developments in this work revolve around a precise ontological fork:

\begin{itemize}[leftmargin=*]
  \item \textbf{HDI / PCC+QHP:} The particle's actual path is a classical stationary-action curve of a bare Lagrangian \(L_c\), selected by a physically meaningful global phase anchored at emission (PCC). The wavefunction (PRQSV) is ontic, evolves through all available paths, and its interference structure modulates which classical path is realized in each run, but the realized path itself is piecewise classical (e.g., straight rays in free space).\,[file:1]

  \item \textbf{de Broglie--Bohm:} The particle's actual path is a continuous streamline of the quantum probability current \(\mathbf j\), equivalently the gradient of the phase \(S\) of \(\psi = R e^{iS/\hbar}\), and is stationary for an \emph{effective} action that includes the quantum potential \(Q\).\,[web:5][web:52][web:69]
  In regions of interference, these trajectories can bend dramatically and are not classical stationary-action paths of \(L_c\).
\end{itemize}

Both theories are deterministic, single-world, and empirically equivalent at the level of \(|\psi|^2\), but they make \emph{different claims about the ontic history} of a particle between preparation and measurement.\,[file:1][web:52][web:71]

\subsection{Main contributions}

The report crystallizes this fork mathematically and numerically:

\begin{enumerate}[leftmargin=*]
  \item \textbf{Path-integral RCC formalization of PCC+QHP:} We reinterpret Gurappa's PCC+QHP in the language of Feynman path integrals: a fixed global phase \(\Phi\) anchors the phase of the initial state, and QHP selects a single path \(r_\Phi(t)\) from the classical stationary set \(\{\delta S_c [r] = 0\}\), where \(S_c=\int L_c\,dt\).\,[file:1][web:31][web:34][web:35]

  \item \textbf{Explicit incompatibility with Bohmian trajectories:} We show that Bohmian trajectories are stationary for an effective action \(S_{\text{eff}}=\int(L_c-Q)\,dt\), not for \(S_c\) alone.\,[web:52][web:69]
  In double-slit and MZI interferometers, Bohmian trajectories bend through interference nodes, whereas PCC+QHP paths are constrained to be classical rays; the disagreement is not just interpretational but dynamical.

  \item \textbf{Unified 2D TDSE simulator architecture:} We present a single numerical framework --- a 2D time-dependent Schr\"odinger solver with split-step Fourier propagation --- that supports both double-slit and MZI geometries, along with:
  \begin{itemize}
    \item Bohmian trajectories from \(\mathbf v=\mathbf j/\rho\), with bilinear interpolation and smooth nodal masking.\,[web:101][web:120][web:123]
    \item A retrodictive PCC+QHP ray-assignment layer that uses a hidden global phase to choose between classical stationary rays per detection event.
    \item A coherence slider representing \(\gamma = \langle m_1|m_2\rangle\) of which-path marker states, implementing decoherence and quantum-eraser scenarios.\,[web:81][web:90][web:156][web:157]
  \end{itemize}

  \item \textbf{Interactive pedagogical design:} We define and implement (in separate code artifacts) an interactive tool where the underlying \(|\psi|^2\) and coherence evolution are fixed, but users can toggle:
  \begin{itemize}
    \item density only;
    \item density plus Bohmian streamlines;
    \item density plus PCC+QHP straight rays;
    \item both trajectory layers simultaneously.
  \end{itemize}
  This makes the ontological divergence visually obvious: the same interference image sits atop incompatible path ontologies.\,[web:67][web:71][web:92]
\end{enumerate}

As such, the combined analytical and numerical design described here should be treated as prior art for any attempt to patent:
(i) PCC+QHP-based trajectory selection on top of a TDSE solver in interferometric setups;
(ii) interactive visualizations that explicitly juxtapose Bohmian and PCC+QHP ontologies using a shared wave evolution and coherence controls;
(iii) specific combinations of hidden global phase selection rules with classical stationary-path ensembles implementing HDI in software.

\section{Mathematical Framework}

\subsection{Standard path integral and stationary phase}

For a scalar non-relativistic particle with classical Lagrangian \(L_c(r,\dot r,t)\), the propagator from \((r_i,t_i)\) to \((r_f,t_f)\) can be written as the path integral
\begin{equation}
K(r_f,t_f;r_i,t_i) = \int \mathcal{D}[r(t)]\; e^{\frac{i}{\hbar} S_c[r]},\quad
S_c[r] = \int_{t_i}^{t_f} L_c(r,\dot r,t)\,dt.\label{eq:path-int}
\end{equation}
In the semiclassical limit, dominant contributions arise from stationary-action paths \(r_{\text{cl}}(t)\) satisfying
\begin{equation}
\delta S_c[r_{\text{cl}}] = 0,\label{eq:classical-action}
\end{equation}
which are the solutions of the Euler--Lagrange equations for \(L_c\).\,[web:31][web:34]
In standard QM, these paths are approximating {\em saddles} in a complex phase sum; no single path is privileged ontically.

\subsection{Phase Consistency Criterion (PCC)}

In HDI, Gurappa introduces the Phase Consistency Criterion as follows:\,[file:1]
\begin{equation}
\mathrm{ph}(\ket{\psi}) = \mathrm{ph}(\bra{r_i,t_i}\psi\rangle),
\end{equation}
where \(\mathrm{ph}(\cdot)\) denotes the global phase and \((r_i,t_i)\) is the emission event.
We can formalize this in path-integral terms by defining:

\begin{itemize}[leftmargin=*]
  \item A fixed initial amplitude
  \(\langle r_i,t_i|\psi\rangle = A_0 e^{i\theta_0}\).
  \item A hidden global phase variable \(\Phi\) constrained by
  \(\Phi \equiv \theta_0 \mod 2\pi\).
  \item For any path \(r(t)\), a total phase
  \(\phi_\Phi[r] = \Phi + S_c[r]/\hbar\).
\end{itemize}

PCC thus \emph{anchors} the origin of all path phases to the physical emission event, treating global phase as a physically meaningful hidden variable rather than a pure gauge.\,[file:1][web:35]

\subsection{Quantum Hamilton's Principle (QHP)}

Gurappa's QHP asserts that the actual path is the one that renders the accumulated quantum phase stationary:\,[file:1]
\begin{equation}
\delta \left(\int_{t_i}^{t_f} L_c(r(t), \dot r(t))\,dt\right) = 0.\label{eq:QHP}
\end{equation}
In the path-integral language, this means that the realized trajectory \(r_\Phi(t)\) lies in the stationary set
\begin{equation}
\mathcal{P}_{\mathrm{stat}} = \{r(t)\in\mathcal{P}\mid \delta S_c[r]=0\}.
\end{equation}
If there are multiple stationary paths consistent with boundary conditions, QHP plus the global phase \(\Phi\) is intended to select one path \(r_\Phi(t)\) as the ontic trajectory.

Formally, we can define a selection functional
\begin{equation}
r_\Phi = \mathcal{S}_\Phi[\mathcal{P}_{\mathrm{stat}}],
\end{equation}
where the precise nature of \(\mathcal{S}_\Phi\) (how \(\Phi\) discriminates among degenerate stationary paths) is left unspecified in the original HDI paper and is instantiated here numerically by a phase-consistency rule at detection.

\subsection{PRQSV and multiplicity of non-events}

HDI treats the position representation of the Heisenberg state vector (PRQSV) as a physically real wavefunction that evolves through all paths while the particle follows a single trajectory inside it.\,[file:1]
In the path-integral view, this means:
\begin{equation}
\psi_\Phi(r_f,t_f) = \int_{\mathcal{P}(r_i\to r_f)} \mathcal{D}[r]\; A[r]\; e^{i\phi_\Phi[r]},
\end{equation}
but only \(r_\Phi(t)\in\mathcal{P}_{\mathrm{stat}}\) is realized as an event trajectory.

All other paths contribute to the PRQSV as \emph{counterfactual non-events}: they shape the interference pattern (through the phase structure of the superposition), but do not correspond to actual spacetime histories of the particle in a given run.\,[file:1]

\subsection{Bohmian trajectories and effective action}

De Broglie--Bohm theory begins from the polar decomposition \(\psi = R e^{iS/\hbar}\), with \(R\ge 0\) and real \(S\).
Writing the Schr\"odinger equation in this form yields a continuity equation and a quantum Hamilton--Jacobi equation:\,[web:5][web:52]
\begin{align}
\frac{\partial R^2}{\partial t} + \nabla\cdot\left(\frac{R^2}{m}\nabla S\right) &= 0,\\
\frac{\partial S}{\partial t} + \frac{(\nabla S)^2}{2m} + V + Q &= 0,
\end{align}
with quantum potential
\begin{equation}
Q(r,t) = -\frac{\hbar^2}{2m} \frac{\nabla^2 R}{R}.
\end{equation}

The Bohmian velocity field is
\begin{equation}
\dot r(t) = \frac{1}{m}\nabla S(r(t),t),
\end{equation}
and can equivalently be written using probability current \(\mathbf j\) and density \(\rho\):
\begin{equation}
\mathbf j = \frac{\hbar}{m}\operatorname{Im}(\psi^* \nabla\psi),\quad
\mathbf v = \frac{\mathbf j}{|\psi|^2}.\,[web:5][web:120][web:123]
\end{equation}

This dynamics can be regarded as the stationary condition of an \emph{effective} action
\begin{equation}
S_{\text{eff}}[r] = \int (L_c - Q)\,dt,
\end{equation}
so that the Bohmian trajectories are not generally stationary for the bare classical action \(S_c\) used in QHP, but for \(S_c - \int Q\,dt\).\,[web:52][web:69][web:108]

\subsection{Core incompatibility}

We can now make the tension precise:

\begin{itemize}[leftmargin=*]
  \item HDI's QHP (as written) uses classical stationarity:
  \[
  \delta S_c[r_\Phi] = 0.
  \]
  The realized paths in free space between optical elements are therefore straight-line rays.

  \item Bohmian trajectories obey stationarity for \(S_{\text{eff}}\), where \(Q\neq 0\) in regions of interference; they are curved streamlines of the quantum current that avoid nodes of \(R\) and can exhibit ``surreal'' behavior.\,[web:71][web:88][web:92]
\end{itemize}

Whenever \(Q\) is appreciable (e.g., in the interference region of a double-slit or MZI), the two stationarity conditions define different Euler--Lagrange problems, and the resulting ontic trajectories differ even while both reproduce the same \(|\psi|^2\) on the screen.\,[web:67][web:71]

In particular, in a double-slit setup:

\begin{itemize}[leftmargin=*]
  \item PCC+QHP demands that the actual path to a point \(x\) on the screen be a straight classical ray from the source through either slit 1 or slit 2.

  \item Bohmian integration of \(\dot{\mathbf r} = \nabla S/m\) yields curves that weave between intensity nodes; some trajectories starting near slit 1 can land in regions geometrically associated with slit 2 and vice versa.\,[web:68][web:71][web:74]
\end{itemize}

The situation is analogous in an MZI: HDI insists on piecewise classical arm trajectories, while Bohmian trajectories can bend in the recombination region.

\section{Interferometric Geometries}

\subsection{Double-slit}

We model a standard double-slit arrangement in a 2D \((x,z)\) plane: slits at \(z = z_s\) with centers at \(x=\pm a/2\), screen at \(z=L\), and a monochromatic beam propagating in \(+z\).

In the paraxial/Fraunhofer regime, the screen wavefunction can be approximated as
\begin{equation}
\psi(x,L) \propto \psi_1(x,L) + \psi_2(x,L),
\end{equation}
with each slit contributing a Gaussian-diffracted term
\begin{equation}
\psi_{1,2}(x,L) \approx A \exp\left[-\frac{(x\mp a/2)^2}{4\sigma^2}\right]
\exp\left[i k L + i\frac{k(x \mp a/2)^2}{2L}\right],
\end{equation}
yielding the interference intensity
\begin{equation}
I(x) = |\psi_1 + \psi_2|^2
\simeq 4C^2\cos^2\left(\frac{k a}{2L}x\right)\cdot f_{\text{env}}(x),\label{eq:double-slit-intensity}
\end{equation}
where \(f_{\text{env}}(x)\) is a slowly varying envelope.\,[web:67][web:73][web:76]

\subsubsection{HDI picture}

\begin{itemize}[leftmargin=*]
  \item PRQSV: the wavefunction \(\psi_\Phi(x,z)\) passes through \emph{both} slits.
  \item QHP: the actual particle trajectory is a classical straight ray from the source through one slit (say slit 2) to \((x,L)\).
  \item The amplitude associated with slit 1 is a counterfactual non-event that nonetheless contributes interference in \(\psi\) and thus modulates the \emph{probability} that the slit-2 ray is realized at that \(x\).\,[file:1]
\end{itemize}

\subsubsection{Bohmian picture}

\begin{itemize}[leftmargin=*]
  \item PRQSV: same as above; the wavefunction obeys the TDSE and produces the usual interference pattern.
  \item The actual trajectory is obtained by integrating the Bohmian velocity:
  \begin{equation}
  \dot{\mathbf r} = \frac{\mathbf j}{|\psi|^2}.\label{eq:bohm-velocity}
  \end{equation}
  \item In the interference region, \(R=|\psi|\) develops nodes and fringes; the resulting streamlines bend accordingly and are not straight classical rays.\,[web:68][web:71][web:74]
\end{itemize}

\subsection{Mach--Zehnder Interferometer (MZI)}

A Mach--Zehnder interferometer splits an incoming beam into two spatially separated arms and recombines them at a second beam splitter.\,[web:39][web:89][web:93]

At the mode level (ignoring transverse structure), one can write after BS1:
\begin{equation}
\ket{\psi_1} = \frac{1}{\sqrt{2}}\big(i\ket{u} + \ket{d}\big),
\end{equation}
with a phase shifter in the upper arm introducing a relative phase \(\phi\):
\begin{equation}
\ket{\psi_2} = \frac{1}{\sqrt{2}}\big(i e^{i\phi}\ket{u} + \ket{d}\big).
\end{equation}
After BS2, the output amplitudes at detectors \(D_1\) (upper) and \(D_2\) (lower) are
\begin{equation}
c_1(\phi) = \frac{i}{2}(1+e^{i\phi}),\quad
c_2(\phi) = \frac{1}{2}(1-e^{i\phi}),
\end{equation}
giving probabilities
\begin{equation}
P_1(\phi) = \cos^2\left(\frac{\phi}{2}\right),\quad
P_2(\phi) = \sin^2\left(\frac{\phi}{2}\right).\,[web:39]
\end{equation}

In a full 2D free-space implementation, the two arms are represented as spatially separated beams recombining at BS2.

\subsubsection{HDI picture}

\begin{itemize}[leftmargin=*]
  \item PRQSV: the wavefunction splits into both arms, accumulating phases and recombining at BS2.
  \item QHP: the actual path is a classical ray running along one arm; no free-space bending in the overlap region beyond what classical optics prescribes.
  \item Interference at BS2 modulates which output port is realized; the unused arm contribution is a counterfactual non-event.
\end{itemize}

\subsubsection{Bohmian picture}

In a spatial representation, as the beams overlap near BS2:

\begin{itemize}[leftmargin=*]
  \item The Bohmian streamlines curve according to the local phase structure, sometimes bending from one arm into the bright output port in ways that are not simply ``which arm'' dependent.\,[web:92]
  \item The sorting into bright vs.\ dark ports is implemented by the \emph{continuous bending} of individual trajectories rather than by a discrete choice among straight arms.
\end{itemize}

Again, the ontological difference is: HDI keeps trajectories piecewise classical and uses global phase selection; Bohm bends trajectories through the quantum phase field.

\section{Unified Numerical Architecture}

\subsection{2D TDSE solver with split-step Fourier method}

We consider a 2D TDSE for a wavefunction \(\psi(x,z,t)\) (propagation along \(z\) treated as ``time'' here for visualization, though in code we use physical time \(t\)):
\begin{equation}
i\hbar\frac{\partial}{\partial t} \psi(x,z,t)
= \left[-\frac{\hbar^2}{2m}(\partial_x^2 + \partial_z^2) + V(x,z)\right]\psi(x,z,t).
\end{equation}

Using the split-step Fourier method, one time step \(\Delta t\) is approximated by:
\begin{align}
\psi'(x,z) &= \exp\left(-\frac{i}{2\hbar}V(x,z)\Delta t\right)\psi(x,z,t),\\
\tilde{\psi}'(k_x,k_z) &= \mathcal{F}[\psi'](k_x,k_z),\\
\tilde{\psi}''(k_x,k_z) &= \exp\left(-\frac{i}{2m\hbar} (k_x^2+k_z^2)\Delta t\right)\tilde{\psi}'(k_x,k_z),\\
\psi''(x,z) &= \mathcal{F}^{-1}[\tilde{\psi}''](x,z),\\
\psi(x,z,t+\Delta t) &= \exp\left(-\frac{i}{2\hbar}V(x,z)\Delta t\right)\psi''(x,z).\label{eq:split-step}
\end{align}
Absorbing boundaries are implemented via a multiplicative mask near the edges.\,[web:79][web:96][web:101][web:105]

\subsection{Potential geometries}

We implement:

\begin{itemize}[leftmargin=*]
  \item \textbf{Double slit:} A barrier at \(z=z_b\) with two slit openings, encoded as a high, finite potential outside the slits:
  \[
  V_{\text{DS}}(x,z) =
  \begin{cases}
    V_0 & \text{for } z_b \le z \le z_b+\Delta z,~x \notin \text{slits},\\
    0 & \text{otherwise.}
  \end{cases}
  \]

  \item \textbf{MZI:} Two barrier planes representing BS1 and BS2 with slit openings for the arms, plus optional phase shifters in one arm.\,[web:93]
\end{itemize}

\subsection{Bohmian trajectories from \(\mathbf j/\rho\)}

At each output frame, we compute the probability density \(\rho = |\psi|^2\) and current components:
\begin{align}
j_x &= \frac{\hbar}{m}\operatorname{Im}\left(\psi^*\partial_x\psi\right),\\
j_z &= \frac{\hbar}{m}\operatorname{Im}\left(\psi^*\partial_z\psi\right),
\end{align}
then obtain the velocity field
\begin{equation}
v_x = \frac{j_x}{\rho},\; v_z = \frac{j_z}{\rho},
\end{equation}
with appropriate regularization for small \(\rho\).\,[web:101][web:120][web:123]

Numerically:

\begin{itemize}[leftmargin=*]
  \item use central differences for \(\partial_x\psi\) and \(\partial_z\psi\);
  \item define \(\rho_{\rm safe} = \max(\rho,\rho_{\rm floor})\) for some small \(\rho_{\rm floor}\);
  \item implement bilinear interpolation of \(v_x,v_z\) to particle positions to avoid grid locking;
  \item smoothly gate velocities near nodes by a factor \(g(\rho)\in[0,1]\) that tends to zero as \(\rho\to 0\).\,[web:126][web:155]
\end{itemize}

Particle positions \((x_p,z_p)\) are updated using an ODE integrator (e.g., Euler or RK4) with time step \(\Delta t_{\rm traj}\le \Delta t\).

\subsection{PCC+QHP retrodictive ray assignment}

After the wave packet hits the screen at \(z=z_f\), we define a screen intensity profile.

For the double slit, we extract approximate slit-resolved contributions \(\psi_1(x)\), \(\psi_2(x)\) and model decoherence via a complex coherence factor \(\gamma\in[0,1]\):
\begin{equation}
I(x) = |\psi_1|^2 + |\psi_2|^2 + 2\,\Re\left(\gamma\,\psi_1^*\psi_2\right).\label{eq:coherence-intensity}
\end{equation}
Sampling detection positions \(x_f\) from \(I(x)\) yields an ensemble of detection events.\,[web:81][web:90][web:156]

For each such event, we assign a PCC+QHP path as follows:

\begin{enumerate}[leftmargin=*]
  \item Compute classical path lengths for paths via slit 1 and slit 2:
  \[
  L_1(x_f) = |r_s - r_{s1}| + |r_{s1}-r_f|,\quad
  L_2(x_f) = |r_s - r_{s2}| + |r_{s2}-r_f|,
  \]
  where \(r_s\) is source, \(r_{s1}, r_{s2}\) are slit positions, \(r_f\) is detection point.

  \item Define classical phases
  \[
  \phi_1(x_f) = k L_1(x_f),\quad \phi_2(x_f) = k L_2(x_f).
  \]

  \item Draw a hidden emission phase \(\Phi\in[0,2\pi)\).

  \item Choose slit index \(j\in\{1,2\}\) that minimises \(|\mathrm{wrap}(\phi_j(x_f)+\Phi)|\), where \(\mathrm{wrap}\) maps angles to \((-\pi,\pi]\).

  \item Assign the realized trajectory as the classical ray:
  \[
  \text{ray}_\Phi(x_f) = r_s \to r_{sj} \to r_f.
  \]
\end{enumerate}

This is a minimal numerical implementation of PCC+QHP: the PRQSV supplies intensities and phases; the hidden global phase \(\Phi\) plus classical action phases \(\phi_j\) select one stationary classical path per detection.

\subsection{Coherence and marker states}

To model which-path marking and quantum erasure, we introduce a two-component wavefunction
\begin{equation}
\Psi(x,z) = \psi_1(x,z)\ket{m_1} + \psi_2(x,z)\ket{m_2},
\end{equation}
with marker states \(\ket{m_1},\ket{m_2}\) in a separate two-dimensional Hilbert space.\,[web:151][web:156][web:157]
The reduced density on the screen after tracing out the marker is
\begin{equation}
I(x) = |\psi_1|^2 + |\psi_2|^2 + 2 \Re\left( \langle m_1|m_2\rangle \psi_1^*\psi_2 \right),
\end{equation}
so we can identify the coherence factor
\(\gamma=\langle m_1|m_2\rangle\) in Eq.~\eqref{eq:coherence-intensity}.

A quantum eraser is implemented by a basis change in the marker space; conditional on eraser outcomes, interference reappears in subensembles even if the unconditional pattern is washed out.\,[web:151][web:156][web:157]

\section{Illustrative Code Snippets}

\subsection{Bohmian velocity and particle advection}

The following Python-style fragment shows how to compute the current, velocity, and advect Bohmian particles on a 2D grid:

\begin{verbatim}
import numpy as np

def prob_current(psi, dx, dz, hbar, m):
    grad_x = (psi[2:,1:-1] - psi[:-2,1:-1]) / (2*dx)
    grad_z = (psi[1:-1,2:] - psi[1:-1,:-2]) / (2*dz)
    jx = np.zeros_like(psi, dtype=float)
    jz = np.zeros_like(psi, dtype=float)
    jx[1:-1,1:-1] = (hbar/m) * np.imag(np.conj(psi[1:-1,1:-1]) * grad_x)
    jz[1:-1,1:-1] = (hbar/m) * np.imag(np.conj(psi[1:-1,1:-1]) * grad_z)
    return jx, jz

def velocity_field(psi, dx, dz, hbar, m, rho_floor=1e-12):
    rho = np.abs(psi)**2
    jx, jz = prob_current(psi, dx, dz, hbar, m)
    rho_safe = np.maximum(rho, rho_floor)
    vx = jx / rho_safe
    vz = jz / rho_safe
    # Smooth gating near nodes
    g = rho_safe / (rho_safe + rho_floor)
    vx *= g; vz *= g
    return vx, vz, rho

def bilinear_interpolate(field, x_coords, z_coords, x_grid, z_grid):
    # x_coords, z_coords: particle positions; x_grid, z_grid: 1D arrays
    Nx, Nz = field.shape
    dx = x_grid[1] - x_grid[0]
    dz = z_grid[1] - z_grid[0]

    ix = (x_coords - x_grid[0]) / dx
    iz = (z_coords - z_grid[0]) / dz

    i0 = np.floor(ix).astype(int); i1 = i0 + 1
    j0 = np.floor(iz).astype(int); j1 = j0 + 1
    i0 = np.clip(i0, 0, Nx-1); i1 = np.clip(i1, 0, Nx-1)
    j0 = np.clip(j0, 0, Nz-1); j1 = np.clip(j1, 0, Nz-1)

    sx = ix - i0; sz = iz - j0
    f00 = field[i0, j0]; f10 = field[i1, j0]
    f01 = field[i0, j1]; f11 = field[i1, j1]

    f0 = (1-sx)*f00 + sx*f10
    f1 = (1-sx)*f01 + sx*f11
    return (1-sz)*f0 + sz*f1

def advect_particles(positions, vx, vz, x_grid, z_grid, dt, rho, rho_cut=1e-8):
    new_pos = []
    xs = np.array([p[0] for p in positions])
    zs = np.array([p[1] for p in positions])

    # interpolate velocity at particle positions
    vx_p = bilinear_interpolate(vx, xs, zs, x_grid, z_grid)
    vz_p = bilinear_interpolate(vz, xs, zs, x_grid, z_grid)
    rho_p = bilinear_interpolate(rho, xs, zs, x_grid, z_grid)

    mask = rho_p > rho_cut
    xs_new = xs + vx_p * dt
    zs_new = zs + vz_p * dt

    for x_new, z_new, keep in zip(xs_new, zs_new, mask):
        if keep:
            new_pos.append((x_new, z_new))
    return new_pos
\end{verbatim}

This realizes Bohmian trajectories numerically in a manner aligned with prior work on TDSE-based trajectory simulations.\,[web:101][web:120][web:123]

\subsection{Retrodictive PCC+QHP path assignment}

The following snippet demonstrates the PCC+QHP ray assignment at the screen:

\begin{verbatim}
import numpy as np

def classical_path_phase(xf, source, slit, screen_z, k):
    # source, slit: (x,z) tuples; xf: screen x coordinate
    xs, zs = source
    xsj, zsj = slit
    L1 = np.hypot(xsj - xs, zsj - zs)
    L2 = np.hypot(xf  - xsj, screen_z - zsj)
    return k * (L1 + L2)

def wrap_angle(theta):
    # Map angle to (-pi, pi]
    return (theta + np.pi) % (2*np.pi) - np.pi

def pcc_assign_rays(x_hits, source, slit1, slit2, screen_z, k):
    rays = []
    Phi = np.random.uniform(0.0, 2*np.pi, size=len(x_hits))
    for xf, phi in zip(x_hits, Phi):
        phi1 = classical_path_phase(xf, source, slit1, screen_z, k)
        phi2 = classical_path_phase(xf, source, slit2, screen_z, k)
        d1 = abs(wrap_angle(phi1 + phi))
        d2 = abs(wrap_angle(phi2 + phi))
        slit = slit1 if d1 <= d2 else slit2
        rays.append({
            "xf": float(xf),
            "slit": 1 if slit is slit1 else 2,
            "ray": [source, slit, (xf, screen_z)]
        })
    return rays
\end{verbatim}

This algorithm operationalizes PCC+QHP as a deterministic hidden-variable selection rule on top of a given wave/field solution.

\section{Prior Art and Defensive Scope}

\subsection{Novel combinations}

The following combinations are, to the best of our knowledge, not standard in prior literature and thus are explicitly disclosed:

\begin{enumerate}[leftmargin=*]
  \item \textbf{A unified 2D interferometer simulator} that:
  \begin{itemize}
    \item serves both double-slit and Mach--Zehnder geometries,
    \item uses a \emph{single} TDSE solver (split-step Fourier or CN) for wave evolution,
    \item overlays de Broglie--Bohm trajectories computed from \(\mathbf j/\rho\),
    \item and overlays PCC+QHP retrodictive classical rays based on a hidden global phase,
    \item all within one interactive visual frame where the user can toggle layers and adjust coherence.
  \end{itemize}

  \item \textbf{A coherence slider interpreted as \(\langle m_1|m_2\rangle\)} of marker states, controlling both the interference term and the distribution over PCC+QHP path selections, while keeping individual rays straight and Bohmian streamlines curved.

  \item \textbf{A specific phase-based ray selection algorithm} that:
  \begin{itemize}
    \item computes classical path phases via slit positions,
    \item samples a hidden global phase \(\Phi\),
    \item and selects the path minimizing \(|\mathrm{wrap}(\phi_j(x_f)+\Phi)|\) per detection.
  \end{itemize}

  \item \textbf{An explicit comparison of HDI vs.\ dBB at the level of variational principles}, highlighting that dBB trajectories are stationary for \(S_c - \int Q\,dt\), not \(S_c\), and using this mismatch to distinguish path ontologies in interference regions.
\end{enumerate}

\subsection{Implications}

Any patent claim attempting to monopolize:

\begin{itemize}[leftmargin=*]
  \item the use of PCC+QHP or HDI-inspired hidden global phase rules for trajectory assignment on top of TDSE-based interferometer simulations,
  \item interactive educational tools that simultaneously visualise Bohmian and HDI/PCC+QHP trajectories over a shared wave evolution with coherence/decoherence controls,
  \item or specific algorithms for retrodictive classical-path assignment conditioned on phases extracted from a quantum field,

would need to contend with this publication as prior art.

\section{Recommendations and Extensions}

\subsection{Time-dependent quantum eraser}

The next logical extension is a time-dependent quantum eraser implemented in the same framework:

\begin{itemize}[leftmargin=*]
  \item Evolve a two-component wavefunction \((\psi_A,\psi_B)\) entangled with path markers.
  \item Implement basis changes in the marker space to produce erasure and conditional interference in subsets of runs.
  \item Track Bohmian trajectories under this extended multi-component field and compare with PCC+QHP rays assigned from the same final density and marker outcomes.\,[web:151][web:156][web:157]
\end{itemize}

\subsection{Entangled pairs and Bell tests}

The framework may be generalized to two-particle entangled interferometers, where HDI's QHP in configuration space would need to specify stationary paths under nonlocal boundary conditions.
Comparing this with Bohmian configuration-space trajectories would probe the feasibility of HDI as a nonlocal hidden-variable model consistent with Bell tests.\,[web:23][web:50][web:54]

\appendix
\section*{Mathematical Appendix}
\addcontentsline{toc}{section}{Mathematical Appendix}

\section{Operator Formulation of PCC+QHP}

\subsection{Global phase anchoring and projectors}

Let \(\mathcal{H}\) be the system Hilbert space and \(U(t,t_0)\) the unitary propagator generated by a self-adjoint Hamiltonian \(H\) on a dense domain \(\mathcal{D}(H)\subset\mathcal{H}\).
We fix an emission event \((r_0,t_0)\) and define the normalized initial state
\[
\ket{\psi_0} = \ket{\psi(t_0)}\in\mathcal{H},\qquad \|\psi_0\|=1.
\]
Let \(\Pi_{r_0}\) be the spectral projector associated with the position observable at \(r_0\), formally \(\Pi_{r_0}=\ket{r_0}\bra{r_0}\) in the usual generalized sense.

The Phase Consistency Criterion (PCC) can be expressed as the constraint
\begin{equation}
\mathrm{ph}(\braket{r_0|\psi_0}) = \mathrm{ph}(\braket{\psi_0|\psi_0}),\label{eq:PCC-operator}
\end{equation}
i.e.\ the global phase of \(\ket{\psi_0}\) is fixed by the complex argument of its amplitude at the emission point.
Operationally, this picks a global phase convention:
\[
\ket{\psi_0^\Phi} := e^{-i\theta_0}\ket{\psi_0},\quad
\theta_0 := \arg\braket{r_0|\psi_0},
\]
so that \(\braket{r_0|\psi_0^\Phi}\in\mathbb{R}_{\ge 0}\).
This defines a one-parameter family of hidden variables \(\Phi\in[0,2\pi)\) via
\[
\ket{\psi_0^\Phi} = e^{i\Phi}\ket{\psi_0^\mathrm{ref}},
\]
for some fixed reference state \(\ket{\psi_0^\mathrm{ref}}\) with \(\braket{r_0|\psi_0^\mathrm{ref}}\in\mathbb{R}_{\ge 0}\).

\begin{lemma}[PCC Anchoring]
For any normalized \(\ket{\psi_0}\) with \(\braket{r_0|\psi_0}\neq 0\), there exists a unique \(\Phi\in[0,2\pi)\) such that
\[
\braket{r_0|\psi_0^\Phi} = |\braket{r_0|\psi_0}|\in\mathbb{R}_{\ge 0},\qquad
\ket{\psi_0^\Phi} = e^{i\Phi}\ket{\psi_0}.
\]
\end{lemma}

\begin{proof}
Take \(\theta_0 = \arg\braket{r_0|\psi_0}\).
Define \(\Phi = -\theta_0\) modulo \(2\pi\) and \(\ket{\psi_0^\Phi}=e^{i\Phi}\ket{\psi_0}\).
Then
\[
\braket{r_0|\psi_0^\Phi} = e^{i\Phi}\braket{r_0|\psi_0}
= e^{-i\theta_0}\braket{r_0|\psi_0}
= |\braket{r_0|\psi_0}|\in\mathbb{R}_{\ge 0}.
\]
Uniqueness follows since if \(\Phi'\) yields the same property we must have
\[
e^{i(\Phi'-\Phi)} \in \mathbb{R}_{>0}\quad\Rightarrow\quad \Phi'-\Phi\in 2\pi\mathbb{Z}.
\]
\end{proof}

\subsection{QHP as a constraint on propagators}

Let \(L_c\) be a classical Lagrangian, and denote by \(\mathcal{P}_{\mathrm{stat}}\) the set of classical stationary paths between two spacetime points.
For each \(r_f\) at time \(t_f\), define the formal propagator decomposition
\[
U(t_f,t_0)\ket{\psi_0^\Phi} = \int_{\mathbb{R}^3} \psi_\Phi(r_f,t_f)\ket{r_f}\,d^3r_f,
\]
with
\[
\psi_\Phi(r_f,t_f) = \sum_{r\in\mathcal{P}(r_0\to r_f)} A_r e^{\frac{i}{\hbar}S_c[r] + i\Phi}.
\]
Quantum Hamilton's Principle (QHP) entails that the actual history of a particle between \(t_0\) and \(t_f\) belongs to
\[
\mathcal{P}_{\mathrm{stat}}(r_0\to r_f) := \{r\in\mathcal{P}(r_0\to r_f)\mid \delta S_c[r]=0\}.
\]
Formally, we can define a trajectory-valued map
\[
\Gamma : [0,2\pi)\times \mathbb{R}^3 \to \mathcal{P}_{\mathrm{stat}},\quad
(\Phi,r_f)\mapsto r_\Phi(\cdot;r_f),
\]
such that for each emission phase \(\Phi\) and detection point \(r_f\), \(\Gamma(\Phi,r_f)\) selects a single stationary-action path.
This can be regarded as a measurable selection from the stationary set provided some regularity; we do not require uniqueness at the classical level.

\section{Effective Action and Bohmian Stationarity}

\subsection{Polar decomposition and quantum potential}

Given \(\psi = R e^{iS/\hbar}\) with \(R\ge0\) and real \(S\), the Schr\"odinger equation on \(\mathbb{R}^d\) with potential \(V\) yields the standard decomposition:
\begin{align}
\frac{\partial R^2}{\partial t} + \nabla\cdot\left(\frac{R^2}{m}\nabla S\right) &= 0,\label{eq:continuity-app}\\
\frac{\partial S}{\partial t} + \frac{(\nabla S)^2}{2m} + V + Q &= 0,\label{eq:HJQ-app}
\end{align}
with quantum potential
\begin{equation}
Q = -\frac{\hbar^2}{2m}\frac{\nabla^2 R}{R},
\end{equation}
whenever \(R>0\).

Define the velocity field
\begin{equation}
\mathbf{v} = \frac{1}{m}\nabla S,\label{eq:v-S}
\end{equation}
and the de Broglie--Bohm trajectories by integrating
\begin{equation}
\dot{\mathbf{r}}(t) = \mathbf{v}(\mathbf{r}(t),t).
\end{equation}

\subsection{Action functional with quantum potential}

Consider the effective Lagrangian
\[
L_{\text{eff}}(r,\dot r,t) := \frac{m}{2}|\dot r|^2 - V(r,t) - Q(r,t),
\]
and the corresponding action
\begin{equation}
S_{\text{eff}}[r] = \int_{t_0}^{t_f} L_{\text{eff}}(r(t),\dot r(t),t)\,dt.\label{eq:Seff}
\end{equation}
Formally, the Euler--Lagrange equation is
\begin{equation}
m\ddot r(t) = -\nabla V(r(t),t) - \nabla Q(r(t),t).\label{eq:EL-eff}
\end{equation}
On the other hand, differentiating \eqref{eq:v-S} along a trajectory gives
\begin{equation}
m\ddot r(t) = m\frac{d}{dt}\mathbf{v}(r(t),t)
= m\left(\partial_t \mathbf{v} + (\mathbf{v}\cdot\nabla)\mathbf{v}\right).
\end{equation}
Standard manipulations on \eqref{eq:HJQ-app} show that Eq.~\eqref{eq:EL-eff} is equivalent to this hydrodynamic acceleration relation.\footnote{This is a classical result in the hydrodynamic formulation of quantum mechanics; see, e.g., \emph{Quantum Dynamics with Bohmian Trajectories}.\,[web:108][web:123]}

Thus, under appropriate regularity assumptions on \(R\) and \(S\), Bohmian trajectories \(r(t)\) are stationary paths for the effective action \(S_{\text{eff}}\), not for the purely classical action
\[
S_c[r] = \int_{t_0}^{t_f}\left(\frac{m}{2}|\dot r|^2 - V(r,t)\right)dt.
\]

\begin{proposition}[Effective Stationarity]
Any sufficiently smooth Bohmian trajectory \(r(t)\) generated by \(\dot r = \nabla S/m\) satisfies the Euler--Lagrange equation for \(L_{\text{eff}} = \frac{m}{2}|\dot r|^2 - V - Q\).
Equivalently,
\[
\delta S_{\text{eff}}[r] = 0,\quad \delta S_c[r]\neq 0\text{ in general when }Q\neq 0.
\]
\end{proposition}

\begin{proof}
The equivalence between \eqref{eq:HJQ-app} and \eqref{eq:EL-eff} is standard: taking the gradient of \eqref{eq:HJQ-app} and using \(\mathbf{v}=\nabla S/m\) yields
\[
m\left(\partial_t \mathbf{v}+(\mathbf{v}\cdot\nabla)\mathbf{v}\right) = -\nabla V - \nabla Q.
\]
Comparing with \eqref{eq:EL-eff} shows that trajectories satisfying \(\dot r=\mathbf{v}(r,t)\) obey the Euler--Lagrange equation for \(L_{\text{eff}}\).
The general lack of stationarity for \(S_c\) in regions where \(Q\neq 0\) follows since \(\nabla Q\) appears explicitly in \eqref{eq:EL-eff}.
\end{proof}

\section{Norm Bounds for Split-Step Propagators}

\subsection{Basic operator splitting}

Let \(H = T + V\) with
\[
T = -\frac{\hbar^2}{2m}\Delta,\quad (V\psi)(x) = V(x)\psi(x),
\]
on a suitable domain.
The exact unitary evolution over \(\Delta t\) is
\[
U(\Delta t) = e^{-iH\Delta t/\hbar}.
\]
The first-order Lie--Trotter splitting uses
\[
U_{\text{LT}}(\Delta t) = e^{-iV\Delta t/\hbar} e^{-iT\Delta t/\hbar},
\]
while the symmetric (Strang) split-step used in the main text is
\[
U_{\text{SS}}(\Delta t) = e^{-iV\Delta t/(2\hbar)} e^{-iT\Delta t/\hbar} e^{-iV\Delta t/(2\hbar)}.\,[web:136][web:138]
\]

\subsection{Operator-norm error bound}

Assume \(T\) and \(V\) are self-adjoint and the commutators \([T,[T,V]]\) and \([V,[T,V]]\) extend to bounded operators on the relevant domain.\footnote{This is a strong assumption, but it is satisfied in many numerically regularized settings; more general treatments exist in the Trotter--Kato theory.\,[web:103][web:121][web:155]}
Then the Baker--Campbell--Hausdorff expansion gives
\[
U(\Delta t) - U_{\text{SS}}(\Delta t)
= \mathcal{O}(\Delta t^3)\quad\text{in operator norm}.
\]
More explicitly, one can write
\begin{equation}
\|U(\Delta t)-U_{\text{SS}}(\Delta t)\|
\le C \Delta t^3\label{eq:SS-bound}
\end{equation}
for some constant \(C\) depending on \(\|[T,[T,V]]\|\) and \(\|[V,[T,V]]\|\).\,[web:136][web:139]

\begin{sketchproof}
Using the formal Strang splitting error expansion
\[
e^{\Delta t(A+B)} - e^{\Delta t A/2}e^{\Delta t B}e^{\Delta t A/2}
= \frac{\Delta t^3}{12}\big([B,[B,A]] + [A,[A,B]]\big) + \mathcal{O}(\Delta t^5),
\]
with \(A=-iT/\hbar\), \(B=-iV/\hbar\), and bounding the commutators in norm yields Eq.~\eqref{eq:SS-bound}.
See e.g.\ standard references on symmetric splitting methods.\,[web:136][web:139][web:155]
\end{sketchproof}

\subsection{Stability and unitarity}

Each factor \(e^{-iV\Delta t/(2\hbar)}\) and \(e^{-iT\Delta t/\hbar}\) is unitary on \(\mathcal{H}\) (or on the discretized Hilbert space in finite-difference/Fourier representations).
Thus
\[
\|U_{\text{SS}}(\Delta t)\|=1,
\]
and the method is norm-conserving at the discrete level, modulo numerical rounding and approximations in the implementation of \(T\) (e.g.\ via FFTs).\,[web:101][web:138]

\section{Bounds for Bohmian Velocity Fields}

\subsection{Current and velocity operators}

On a discrete grid with spacing \(\Delta x\) and \(\Delta z\), we approximate
\[
\partial_x\psi(x_i,z_j) \approx \frac{\psi_{i+1,j}-\psi_{i-1,j}}{2\Delta x},\qquad
\partial_z\psi(x_i,z_j) \approx \frac{\psi_{i,j+1}-\psi_{i,j-1}}{2\Delta z},
\]
and compute discrete currents
\begin{align}
j_x^{(h)} &= \frac{\hbar}{m}\Im\big(\psi^* \partial_x^{(h)}\psi\big),\\
j_z^{(h)} &= \frac{\hbar}{m}\Im\big(\psi^* \partial_z^{(h)}\psi\big).
\end{align}
Let \(\rho = |\psi|^2\) and define a floor \(\rho_{\mathrm{floor}}>0\); then
\[
\rho_{\mathrm{safe}} := \max(\rho,\rho_{\mathrm{floor}}),
\]
and velocities
\[
v_x^{(h)} = \frac{j_x^{(h)}}{\rho_{\mathrm{safe}}},\quad
v_z^{(h)} = \frac{j_z^{(h)}}{\rho_{\mathrm{safe}}}.
\]

\subsection{Uniform bounds under spectral cutoff}

Suppose \(\psi\) is band-limited in Fourier space with support \(|k_x|\le K_x\), \(|k_z|\le K_z\).
Then we have the Bernstein-type bounds
\[
\|\partial_x\psi\|_\infty \le K_x \|\psi\|_\infty,\quad
\|\partial_z\psi\|_\infty \le K_z \|\psi\|_\infty.
\]
Assume \(\|\psi\|_\infty \le M\) and \(\rho_{\mathrm{safe}}\ge \rho_{\mathrm{floor}} >0\).
Then
\begin{align}
|j_x^{(h)}| &\le \frac{\hbar}{m}|\psi||\partial_x\psi|
\le \frac{\hbar}{m} M \cdot K_x M = \frac{\hbar}{m}K_x M^2,\\
|v_x^{(h)}| &= \frac{|j_x^{(h)}|}{\rho_{\mathrm{safe}}}
\le \frac{\hbar}{m}K_x M^2 \cdot \frac{1}{\rho_{\mathrm{floor}}}.
\end{align}
Similarly for \(v_z^{(h)}\) with \(K_z\).

\begin{proposition}[Velocity bound under spectral cutoff]
Under the assumptions above,
\[
\|v_x^{(h)}\|_\infty \le C_x,\quad \|v_z^{(h)}\|_\infty \le C_z,
\]
with
\[
C_x = \frac{\hbar}{m}\frac{K_x M^2}{\rho_{\mathrm{floor}}},\quad
C_z = \frac{\hbar}{m}\frac{K_z M^2}{\rho_{\mathrm{floor}}}.
\]
\end{proposition}

Such bounds ensure that explicit time-stepping schemes for Bohmian particle advection can be made stable by taking time steps satisfying a CFL-like condition:
\[
\Delta t_{\mathrm{traj}} \le \min\left(\frac{\Delta x}{C_x},\frac{\Delta z}{C_z}\right),
\]
modulo interpolation errors.\,[web:101][web:120][web:123]

\section{Bounds for PCC-Based Selection}

\subsection{Phase-difference distributions}

For a given detection position \(x_f\), define classical phase differences
\[
\Delta\phi(x_f) := \phi_1(x_f) - \phi_2(x_f),
\]
where \(\phi_j(x_f) = kL_j(x_f)\) are the classical path phases via slit \(j\).
Let \(\Phi\) be uniform on \([0,2\pi)\).
We consider the random variables
\[
D_j(x_f) := \big|\mathrm{wrap}(\phi_j(x_f) + \Phi)\big|\in[0,\pi].
\]

\begin{lemma}[Symmetry]
If \(\Phi\) is uniform on \([0,2\pi)\), then for fixed \(x_f\) the variables \(D_1(x_f)\) and \(D_2(x_f)\) have the same marginal distribution whenever \(\Delta\phi(x_f)\equiv 0 \mod 2\pi\).
\end{lemma}

\begin{proof}
If \(\Delta\phi(x_f)=2\pi n\) for some integer \(n\), then \(\phi_1(x_f)=\phi_2(x_f)+2\pi n\), and
\[
D_1(x_f) = \big|\mathrm{wrap}(\phi_2(x_f)+2\pi n+\Phi)\big|
= \big|\mathrm{wrap}(\phi_2(x_f)+\Phi)\big| = D_2(x_f),
\]
pointwise in \(\Phi\).
\end{proof}

In general, when \(\Delta\phi(x_f)\neq 0\mod 2\pi\), the event \(\{D_1\le D_2\}\) has probability
\[
P_1(x_f) := \mathbb{P}[D_1(x_f)\le D_2(x_f)] = \frac{1}{2} + \delta(x_f)
\]
for some bias \(\delta(x_f)\) depending continuously on \(\Delta\phi(x_f)\).
While the exact form of \(\delta\) is not particularly illuminating in closed-form, we can bound it uniformly.

\begin{proposition}[Bound on PCC bias]
For each \(x_f\), we have
\[
\left|P_1(x_f) - \frac{1}{2}\right| \le \frac{|\Delta\phi(x_f)|}{2\pi}
\]
when \(|\Delta\phi(x_f)|\le \pi\).
\end{proposition}

\begin{proof}[Sketch]
Consider \(\theta_j(\Phi) := \mathrm{wrap}(\phi_j(x_f)+\Phi)\).
As \(\Phi\) varies uniformly over \([0,2\pi)\), the difference \(\theta_1-\theta_2\) sweeps an interval of length \(|\Delta\phi(x_f)|\) modulo \(2\pi\); the region in \(\Phi\)-space where \(D_1\le D_2\) deviates from symmetry by at most this length, which yields a total measure difference bounded by \(|\Delta\phi|/(2\pi)\).
\end{proof}

Thus, as \(|\Delta\phi(x_f)|\to 0\), the PCC selection rule becomes asymptotically unbiased between the two slits, while for large phase differences it can strongly favor the slit with smaller wrapped phase.
This provides a quantitative control over how sensitively the PCC rule depends on classical path phase differences at a fixed detection point.

\bigskip

\noindent
\textbf{Remark.} A more refined analysis would relate the PCC selection probabilities to the quantum amplitudes \(\psi_j(x_f)\) and show how, under suitable distributions of \(\Phi\), the ensemble of realized paths reproduces the Born rule for \(|\psi_1+\psi_2|^2\); we leave this as a direction for further rigorous work.

\nocite{*}
\bibliographystyle{unsrt}  
\bibliography{references}  


\end{document}
